\section{Attention, context, and in-context learning}
\label{sec:sequence}

Two claims about attention are cited far more often than they are tested: that
attention switches between a \emph{positional} and a \emph{semantic} mechanism,
and that in-context learning \emph{emerges} above a threshold of pretraining
task diversity. Both are usually supported by a single curve at a single model
size, which cannot distinguish a transition from a crossover. This block puts
each on a fine control grid with a real scale path, and adds two measurements a
risk curve cannot make: a causal ablation that separates the two attention
routes, and an attention-entropy comparison against the uniform-attention null.

\subsection{The positional--semantic switch}

\manuscriptfigure{figures/family_positional_semantic.pdf}{%
  \textbf{Which route attention learns.} The order parameter is
  $(S-P)/(S+P)$, where $S$ and $P$ are the variance explained when the
  \emph{other} cue has been destroyed by intervention. It runs from $-1$
  (purely positional) to $+1$ (purely semantic).%
}{fig:positional-semantic}

The usual version of this experiment makes the semantic cue perfect, in which
case it wins for every $p>0$ and the ``transition'' is an artefact of the data
rather than a property of the model. We therefore give \emph{both} cues a
finite reliability: the answer token sits at position zero with probability
$1-p$, and carries its true tag with probability $1-q$, with $q$ fixed. A model
that simply follows the more reliable cue should switch route near $p=q$, so
the measurement has a predicted location to agree or disagree with.

\paragraph{Observation.}
Over \AtlasPositionalSemanticConditions{} conditions at
\AtlasPositionalSemanticSizes{} widths, the mechanism order parameter runs the
full range: it sits at $-1$ when the positional cue is reliable and saturates at
$+1$ once that cue is worse than the semantic one. The sign change is
remarkably stable across widths --- $p=0.235$, $0.239$, $0.247$ and $0.245$ at
widths $24$ through $64$ --- against the registered prediction $p=q=0.30$. The
disorder susceptibility peaks earlier, near $p\approx0.13$, but at three of the
four widths that maximum is a single-point spike and is rejected, so no
finite-size signature survives. The order parameter spans
\AtlasPositionalSemanticSemanticRatio{} times its $95\%$ interval; the verdict
is \emph{\AtlasPositionalSemanticVerdict{}}, grade
\AtlasPositionalSemanticGrade{}.

\paragraph{Interpretation.}
The switch is real, it is measured at the level of mechanism rather than
performance, and it lands close to where a reliability-following model would
put it. The causal ablations are what make this a mechanism claim: the model's
risk barely changes across much of the range, so a risk curve alone would have
shown almost nothing.

The offset between the two locations is suggestive rather than established.
The fluctuation maximum sits below the mean sign change, which is what one
expects when runs begin to \emph{disagree} about which route to take before the
average route has switched. But that maximum is exactly the regime in
which a single excursion dominates the variance, and the spike test rejects it
at three widths out of four. The stable quantity here is the mean
sign change; the fluctuation story needs more seeds before it can be told.

\paragraph{Not established.}
Whether the switch is a transition. No admissible finite-size signature
survives, so the verdict is a response shift: the mechanism reversal is
unambiguous, its criticality is not tested.

\paragraph{Next falsifier.}
Vary $q$ and check that the crossover tracks it. If the sign change stays near
$p\approx0.24$ when $q$ is moved to $0.1$ or $0.5$, the model is not following
cue reliability and the interpretation above is wrong.

\subsection{In-context learning versus task diversity}

A small causal transformer is pretrained on in-context linear regression whose
task vectors are drawn from a pool of size $M$. The order parameter is not the
risk but the normalised alignment between the model's predictions on
\emph{fresh} tasks and the exact in-context ridge predictor computed from the
context --- that is, whether the model has reproduced the algorithm, not
whether it scores well.

\manuscriptfigure{figures/family_icl_task_diversity.pdf}{%
  \textbf{From memorising a task pool to running an algorithm.} The order
  parameter is alignment with the exact in-context ridge predictor on
  \emph{fresh} tasks --- zero for a model that has memorised its pretraining
  pool, one for a model that has reproduced the Bayes-optimal rule.%
}{fig:icl}

\paragraph{Observation.}
Across \AtlasIclTaskDiversityConditions{} conditions the two quantities move
together and in opposite directions. At a pool of one task the model memorises
completely: the fresh-versus-pool risk gap is about $+2$ and the ridge
alignment is indistinguishable from zero. Both change through
$M\approx8$--$64$; the gap reaches zero near $M\approx10^2$ and the alignment
saturates around $0.7$ beyond it, with no further improvement out to an
unbounded pool. The order parameter spans
\AtlasIclTaskDiversitySemanticRatio{} times its $95\%$ interval. The verdict is
\emph{\AtlasIclTaskDiversityVerdict{}}, grade
\AtlasIclTaskDiversityGrade{}.

The disorder susceptibility of the alignment has an interior maximum near
$M\approx4$--$8$ at two of the three widths; at the third the global maximum is
an isolated spike and is rejected, leaving too few sizes to fit a sharpening
exponent.

\paragraph{Interpretation.}
The memorisation-to-algorithm change is visible in a quantity that separates
the two explanations, which a learning curve does not. Alignment with the
Bayes-optimal in-context rule is the right order parameter for this claim
precisely because a memorising model can have low risk on its own pool while
having no alignment on fresh tasks.

The susceptibility maxima that do survive sit at a pool of four to eight tasks,
inside the range usually quoted for the onset of in-context learning --- where
independently seeded models disagree most about whether to memorise or to
generalise. With only two admissible sizes this is a coincidence worth
following up, not a finite-size result.

\paragraph{Not established.}
That the change is an emergence threshold. At the model scale used here the
transformer remains far from the ridge baseline in absolute risk even at an
unbounded pool, so what is measured is the onset of algorithmic behaviour, not
its completion.

\paragraph{Next falsifier.}
Train an order of magnitude longer at fixed $M$. If the alignment curve shifts
left with training budget, what looks like a task-diversity threshold is partly
an optimisation threshold, and the control variable is not the one the claim
names.

\subsection{Induction-head formation}

\paragraph{Observation.}
On the canonical induction task --- predict the token that followed the
previous occurrence of the current token --- a two-layer causal model reaches an
accuracy well above the unigram null across
\AtlasInductionHeadConditions{} conditions, and the second layer's attention at
the query position concentrates on exactly the token after the earlier
occurrence. Both the behavioural lift and the attention-level induction score
fall as the vocabulary grows. The verdict is
\emph{\AtlasInductionHeadVerdict{}}, grade \AtlasInductionHeadGrade{}.

\paragraph{Interpretation.}
The circuit is measured twice, once behaviourally and once by reading the
attention pattern, and the two agree: where the lift is large the attention
score is above its uniform null, and where the lift collapses the attention
score returns to null. That agreement is the evidence that the behavioural
number reflects an induction mechanism rather than some other shortcut.

\paragraph{Not established.}
Abruptness as a critical phenomenon. We record the formation step and the
steepest per-checkpoint gain, and censor runs that never cross the registered
lift threshold, but a sharp rise in training time is not by itself a phase
transition in any control parameter.

\subsection{Does attention select, or dilute?}

\paragraph{Observation.}
Sweeping context length over
$L\in\AtlasContextScalingControlRange{}$, the entropy of the query position's
attention distribution stays measurably below the uniform null $\log L$; the
gap is positive throughout and spans
\AtlasContextScalingSemanticRatio{} times its interval across
\AtlasContextScalingConditions{} conditions. The verdict is
\emph{\AtlasContextScalingVerdict{}}, grade \AtlasContextScalingGrade{}.

\paragraph{Interpretation.}
Attention is genuinely selecting rather than spreading: the entropy does grow
with context, but more slowly than the uniform baseline, so the effective
number of attended tokens grows sub-linearly. Reporting the null alongside the
entropy is the whole point --- an entropy that rises with $L$ is compatible with
no selection at all, and only the gap distinguishes the two.

\paragraph{Next falsifier.}
Extend the context grid by another octave. If the gap stops growing, the
effective attended set has saturated, which is a concrete and testable
statement about how far a fixed-width model can use its context.
